\tcbset{
	theoremstyle/.style={
			enhanced,
			breakable,
			colback=white,
			colframe=white,
			boxrule=0pt,
			borderline west={1.5pt}{0pt}{#1},
			left=9pt,
			right=0pt,
			top=2pt,
			bottom=2pt,
			sharp corners,
			fonttitle=\bfseries\color{#1},
			coltitle=#1,
			attach title to upper={\quad},
			separator sign={}
		}
}


\newtcbtheorem[number within=section]
{thm}{Théorème}
{theoremstyle=TheoremColor}
{thm}


\newtcbtheorem[number within=section]
{cor}{Corollaire}
{theoremstyle=CorColor}
{cor}


\newtcbtheorem[number within=section]
{prop}{Proposition}
{theoremstyle=PropColor}
{prop}


\newtcbtheorem[number within=section]
{lem}{Lemme}
{theoremstyle=LemmaColor}
{lem}


\newtcbtheorem[number within=section]
{defn}{Définition}
{theoremstyle=DefColor}
{defn}


\newtcbtheorem[number within=section]
{ex}{Exemple}
{theoremstyle=ExampleColor}
{ex}


\newtcbtheorem[number within=section]
{alg}{Algorithme}
{theoremstyle=AlgColor}
{alg}


\newtcolorbox{important}{
	enhanced,
	breakable,
	colback=white,
	colframe=white,
	boxrule=0pt,
	borderline west={1.5pt}{0pt}{myblue},
	left=10pt,
	right=0pt,
	sharp corners
}


% Plain theorems

\newtheorem{rem}{Remarque}[section]

\newtheorem{exr}{Exercice}[section]
